\documentclass[10pt,a4paper]{article}
\usepackage{amsmath,amssymb,amsthm,graphicx}
\usepackage{amsmath,amsfonts,amssymb,amsthm,enumerate,hyperref,ytableau,authblk,listings}
\usepackage[capitalize]{cleveref}
\newtheorem{theorem}{Theorem}[section]
\newtheorem{lemma}[theorem]{Lemma}
\newtheorem{corollary}[theorem]{Corollary}
\newtheorem{conjecture}[theorem]{Conjecture}

\theoremstyle{definition}

\newtheorem{definition}[theorem]{Definition}
\providecommand{\coset}[2]{[#1 \mapsto #2]}
\providecommand{\dcoset}[2]{[#1 \leftrightarrow #2]}
\providecommand{\id}{\mathsf{id}}
\providecommand{\cM}{\mathcal{M}}
\title{A degree-two symmetric function on the symmetric group has a certificate of size at most n minus two}
\author{Gilad Chase, Neta Dafni, Yuval Filmus, Nathan Lindzey}
\date{arxiv:2210.00245v2; Creative Commons Attribution 4.0 International}
\begin{document}
\maketitle
\noindent\textit{Editorial note.} The counted statement is the author's labelled environment at source0.tex lines 265-267 and its proof environment opens at line 268; the pairing is by adjacency and the retained proof range is the whole environment.\par\medskip
\noindent\textit{Editorial note.} Preamble: the author's own package lines, theorem-environment family and macro block (source lines 5-29).\par\medskip
\noindent\textit{Editorial note.} Bibliography: the article ships its own compiled bibliography file (main.bbl, a complete thebibliography with 29 entries) and the retained proofs cite keys from it, so the card bundles that file unchanged; no bibliographic entry is edited or re-rendered.\par\medskip
\noindent\textit{Editorial note.} Retained-licence form: this source's licence evidence carries no acceptance record and none may be written on anyone's behalf, so the card points at a task-directory evidence copy in the team-wide four-key reuse form whose basis states that the source's staged record is reused as a material basis only. No status key and no acceptance record are written.\par\medskip
\noindent\textit{Editorial note.} Metadata provenance: source\_version is the value recorded in the source's licence evidence (2210.00245v2); the author list is transcribed from the authors\_index key of the source record of the pinned copy, and the title from its title key.\par\medskip
\noindent\textit{Editorial note.} The source carries no label-number dump and no counter values are imposed, so the printed numbering is the wrapper's default; the identity of the counted item is its author label.\par\medskip
\noindent\textit{Editorial note.} r002 note: r001 failed with an undefined reference raised from inside the counted proof; r002 adds exactly the environment that the citation names, as a context span with its provenance, and changes nothing else. The article's own compiled bibliography (main.bbl) was already bundled and is unrelated to this failure.\par\medskip
\noindent\textit{Editorial scope.} One original result of the pinned author manuscript, counted once: the certificate bound for degree-two functions on the symmetric group (author label lem:C-bound-sym).\par\medskip
\noindent\textit{Allowed standard prerequisites (editorial).}\begin{itemize}
\item Degree of a Boolean function on the symmetric group. The degree of f as a multilinear polynomial in the variables recording the images of the points.. It is the hypothesis bounded by two.
\item Certificate size C(f). The largest size of a common certificate over the members of the family, a certificate being a set of prescribed images.. It is the quantity bounded by n minus two.
\item Influences and sensitivity. The degree bound restricts how the value of f can depend on the variables, so a function of high sensitivity cannot have degree two.. It is the contradiction that closes the argument.
\end{itemize}\medskip
\subsection{Code for Lemma \ref{lem:C-bound-sym}} \label{apx:C-bound-sym}

\begin{lemma} \label{lem:C-bound-sym}
If $n \geq 8$ and $f\colon S_n \to \{0,1\}$ has degree at most~$2$, then $C(f) \leq n-2$.
\end{lemma}

\begin{proof}
The proof is by contradiction. Suppose that $n \geq 8$, that $f\colon S_n \to \{0,1\}$ has degree at most~$2$, and that $C(f) = n-1$. Without loss of generality, $C(f,\id) = n-1$ and $f(\id) = 0$. Then $f((i\;j)) = 1$  for all $i \neq j \in [n]$ (here $(i\;j)$ is the transposition switching $i$ and $j$), since otherwise $\{(k,k) : k \neq i,j\}$ would be a certificate for $\id$ of size $n-2$.

We construct a function $g\colon \{0,1\}^4 \to \{0,1\}$ as follows:
\[
 g(y_1,y_2,y_3,y_4) = f\bigl((1\;2)^{y_1} (3\;4)^{y_2} (5\;6)^{y_3} (7\;8)^{y_4}\bigr).
\]
Here $(i\;j)^0 = \id$, $(i\;j)^1 = (i\;j)$, and the input to $f$ is the product of four powers of this kind.

Since $f$ has degree at most~$2$, it can be written as a polynomial of degree at most~$2$ in the variables $x_{ij}$. Given $y_1,y_2,y_3,y_4$, the value of the variables $x_{ij}$ is:
\begin{align*}
    &x_{11} = x_{22} = 1-y_1 && x_{12} = x_{21} = y_1 \\
    &x_{33} = x_{44} = 1-y_2 && x_{34} = x_{43} = y_2 \\
    &x_{55} = x_{66} = 1-y_3 && x_{56} = x_{65} = y_3 \\
    &x_{77} = x_{88} = 1-y_4 && x_{78} = x_{87} = y_4 
\end{align*}
All other variables are assigned zero. This shows that $g$ can be expressed as a polynomial of degree at most~$2$ in the variables $y_1,y_2,y_3,y_4$. We say that $g$ has \emph{degree at most~$2$}. 

By construction, the function $g$ satisfies the following constraints:
\begin{align*}
    &g(0,0,0,0) = f(\id) = 0 \\
    &g(1,0,0,0) = f((1\;2)) = 1 \\
    &g(0,1,0,0) = f((3\;4)) = 1 \\
    &g(0,0,1,0) = f((5\;6)) = 1 \\
    &g(0,0,0,1) = f((7\;8)) = 1 
\end{align*}
We say that $g$ has \emph{sensitivity $4$}.

There are only $2^{16}$ many functions from $\{0,1\}^4$ to $\{0,1\}$. For each function, we can check whether it has degree at most~$2$ by solving linear equations, since each function has a unique representation as a multilinear polynomial in $y_1,y_2,y_3,y_4$ (this is a basic fact in Boolean function analysis). Going over all such functions in SAGE~\cite{sagemath}, we find out that none of them has sensitivity~$4$, a contradiction. The relevant code appears in \Cref{apx:C-bound-sym}. 
\end{proof}

\newcommand{\etalchar}[1]{$^{#1}$}
\begin{thebibliography}{BMRR21}

\bibitem[AK97]{AK1}
Rudolf Ahlswede and Levon~H. Khachatrian.
\newblock The complete intersection theorem for systems of finite sets.
\newblock {\em European J. Combin.}, 18(2):125--136, 1997.

\bibitem[AK99]{AK2}
Rudolf Ahlswede and Levon~H. Khachatrian.
\newblock A pushing-pulling method: new proofs of intersection theorems.
\newblock {\em Combinatorica}, 19(1):1--15, 1999.

\bibitem[BdW02]{BuhrmandeWolf}
Harry Buhrman and Ronald de~Wolf.
\newblock Complexity measures and decision tree complexity: a survey.
\newblock {\em Theoret. Comput. Sci.}, 288(1):21--43, 2002.
\newblock Complexity and logic (Vienna, 1998).

\bibitem[BMRR21]{BMRR}
Angelot Behajaina, Roghayeh Maleki, Aina~Toky Rasoamanana, and A.~Sarobidy
  Razafimahatratra.
\newblock 3-setwise intersecting families of the symmetric group.
\newblock {\em Discrete Math.}, 344(8):Paper No. 112467, 15, 2021.

\bibitem[CK03]{CameronKu}
Peter~J. Cameron and C.~Y. Ku.
\newblock Intersecting families of permutations.
\newblock {\em European J. Combin.}, 24(7):881--890, 2003.

\bibitem[CSST08]{CSST08}
Tullio Ceccherini-Silberstein, Fabio Scarabotti, and Filippo Tolli.
\newblock {\em Harmonic analysis on finite groups}, volume 108 of {\em
  Cambridge Studies in Advanced Mathematics}.
\newblock Cambridge University Press, Cambridge, 2008.
\newblock Representation theory, Gelfand pairs and Markov chains.

\bibitem[Daf22]{DafniThesis}
Neta Dafni.
\newblock Complexity measures on the symmetric group and beyond.
\newblock Master's thesis, Technion --- Israel Institute of Technology, 2022.

\bibitem[DFL{\etalchar{+}}21]{DFLLV}
Neta Dafni, Yuval Filmus, Noam Lifshitz, Nathan Lindzey, and Marc Vinyals.
\newblock Complexity measures on the symmetric group and beyond.
\newblock In {\em 12th {I}nnovations in {T}heoretical {C}omputer {S}cience
  {C}onference}, volume 185 of {\em LIPIcs. Leibniz Int. Proc. Inform.}, pages
  Art. No. 87, 5. Schloss Dagstuhl. Leibniz-Zent. Inform., Wadern, 2021.

\bibitem[EFP11]{EFP}
David Ellis, Ehud Friedgut, and Haran Pilpel.
\newblock Intersecting families of permutations.
\newblock {\em J. Amer. Math. Soc.}, 24(3):649--682, 2011.

\bibitem[EKR61]{EKR}
P.~Erd\H{o}s, Chao Ko, and R.~Rado.
\newblock Intersection theorems for systems of finite sets.
\newblock {\em Quart. J. Math. Oxford Ser. (2)}, 12:313--320, 1961.

\bibitem[Ell12a]{EllisCameronKu}
David Ellis.
\newblock A proof of the {C}ameron-{K}u conjecture.
\newblock {\em J. Lond. Math. Soc. (2)}, 85(1):165--190, 2012.

\bibitem[Ell12b]{EllisSetwise}
David Ellis.
\newblock Setwise intersecting families of permutations.
\newblock {\em J. Combin. Theory Ser. A}, 119(4):825--849, 2012.

\bibitem[FD77]{DezaFrankl}
P\'{e}ter Frankl and Mikhail Deza.
\newblock On the maximum number of permutations with given maximal or minimal
  distance.
\newblock {\em J. Combinatorial Theory Ser. A}, 22(3):352--360, 1977.

\bibitem[FHKL16]{FHKL}
Yuval Filmus, Hamed Hatami, Nathan Keller, and Noam Lifshitz.
\newblock On the sum of the {$L_1$} influences of bounded functions.
\newblock {\em Israel J. Math.}, 214(1):167--192, 2016.

\bibitem[Fil17]{FilmusComment}
Yuval Filmus.
\newblock A comment on intersecting families of permutations.
\newblock arXiv:1706.10146, 2017.

\bibitem[FMS21]{FMS}
Shaun Fallat, Karen Meagher, and Mahsa~N. Shirazi.
\newblock The {E}rd{\H{o}}s-{K}o-{R}ado theorem for 2-intersecting families of
  perfect matchings.
\newblock {\em Algebr. Comb.}, 4(4):575--598, 2021.

\bibitem[GM09]{GodsilMeagher}
Chris Godsil and Karen Meagher.
\newblock A new proof of the {E}rd{\H{o}}s-{K}o-{R}ado theorem for intersecting
  families of permutations.
\newblock {\em European J. Combin.}, 30(2):404--414, 2009.

\bibitem[GM16]{GodsilMeagherBook}
Chris Godsil and Karen Meagher.
\newblock {\em Erd\H{o}s-{K}o-{R}ado theorems: algebraic approaches}, volume
  149 of {\em Cambridge Studies in Advanced Mathematics}.
\newblock Cambridge University Press, Cambridge, 2016.

\bibitem[GM17]{GodsilMeagherPM}
Chris Godsil and Karen Meagher.
\newblock An algebraic proof of the {E}rd{\H{o}}s-{K}o-{R}ado theorem for
  intersecting families of perfect matchings.
\newblock {\em Ars Math. Contemp.}, 12(2):205--217, 2017.

\bibitem[Lin17]{Lindzey17}
Nathan Lindzey.
\newblock Erd{\H{o}}s-{K}o-{R}ado for perfect matchings.
\newblock {\em European J. Combin.}, 65:130--142, 2017.

\bibitem[Lin18a]{Lindzey18}
Nathan Lindzey.
\newblock Intersecting families of permutations of perfect matchints.
\newblock arXiv:1811.06160, 2018.

\bibitem[Lin18b]{LindzeyThesis}
Nathan Lindzey.
\newblock {\em Matchings and Representation Theory}.
\newblock PhD thesis, University of Waterloo, 2018.

\bibitem[LM04]{LaroseMalvenuto}
Benoit Larose and Claudia Malvenuto.
\newblock Stable sets of maximal size in {K}neser-type graphs.
\newblock {\em European J. Combin.}, 25(5):657--673, 2004.

\bibitem[MM05]{MeagherMoura}
Karen Meagher and Lucia Moura.
\newblock Erd{\H{o}}s-{K}o-{R}ado theorems for uniform set-partition systems.
\newblock {\em Electron. J. Combin.}, 12:Research Paper 40, 12, 2005.

\bibitem[MR21]{MR21}
Karen Meagher and A.~Sarobidy Razafimahatratra.
\newblock The {E}rd{\H{o}}s-{K}o-{R}ado theorem for 2-pointwise and 2-setwise
  intersecting permutations.
\newblock {\em Electron. J. Combin.}, 28(4):Paper No. 4.10, 21, 2021.

\bibitem[N{\"O}03]{Cliquer}
Sampo Niskanen and Patric R.~J. {\"O}sterg{\aa}rd.
\newblock Cliquer user's guide, version 1.0.
\newblock Technical Report T48, Communications Laboratory, Helsinki University
  of Technology, Espoo, Finland, 2003.

\bibitem[O'D14]{ODonnell}
Ryan O'Donnell.
\newblock {\em Analysis of {B}oolean functions}.
\newblock Cambridge University Press, New York, 2014.

\bibitem[Sag01]{Sagan}
Bruce~E. Sagan.
\newblock {\em The symmetric group}, volume 203 of {\em Graduate Texts in
  Mathematics}.
\newblock Springer-Verlag, New York, second edition, 2001.
\newblock Representations, combinatorial algorithms, and symmetric functions.

\bibitem[{Sag}22]{sagemath}
{Sage Developers}.
\newblock {\em {S}ageMath, the {S}age {M}athematics {S}oftware {S}ystem
  ({V}ersion 9.5)}, 2022.
\newblock {\tt https://www.sagemath.org}.

\end{thebibliography}
\end{document}
